\section{Átomos y reacciones químicas: colección de ejercicios (3.º y 4.º ESO)}

11 de octubre de 2026 · @Quique

\subsection{Reglas y datos}

Treinta ejercicios de cálculo con átomos y reacciones: del 1 al 20 para 3.º ESO y del 21 al 30 para 4.º ESO. Cada apartado tiene una única respuesta y está en las soluciones del final. Los datos salen de esta sección; no hace falta nada más. Los decimales se escriben con punto.

\subsubsection*{Elementos que se usan}

\begin{center}
\begin{tabular}{rllr}
\hline
$Z$ & Símbolo & Nombre & Masa atómica ($\mathrm{u}$) \\
\hline
1  & $\mathrm{H}$  & hidrógeno & 1 \\
2  & $\mathrm{He}$ & helio & 4 \\
3  & $\mathrm{Li}$ & litio & 7 \\
5  & $\mathrm{B}$  & boro & 10.8 \\
6  & $\mathrm{C}$  & carbono & 12 \\
7  & $\mathrm{N}$  & nitrógeno & 14 \\
8  & $\mathrm{O}$  & oxígeno & 16 \\
9  & $\mathrm{F}$  & flúor & 19 \\
10 & $\mathrm{Ne}$ & neón & 20 \\
11 & $\mathrm{Na}$ & sodio & 23 \\
12 & $\mathrm{Mg}$ & magnesio & 24 \\
13 & $\mathrm{Al}$ & aluminio & 27 \\
14 & $\mathrm{Si}$ & silicio & 28 \\
15 & $\mathrm{P}$  & fósforo & 31 \\
16 & $\mathrm{S}$  & azufre & 32 \\
17 & $\mathrm{Cl}$ & cloro & 35.5 \\
18 & $\mathrm{Ar}$ & argón & 40 \\
19 & $\mathrm{K}$  & potasio & 39 \\
20 & $\mathrm{Ca}$ & calcio & 40 \\
25 & $\mathrm{Mn}$ & manganeso & 55 \\
26 & $\mathrm{Fe}$ & hierro & 56 \\
29 & $\mathrm{Cu}$ & cobre & 63.5 \\
30 & $\mathrm{Zn}$ & cinc & 65.4 \\
35 & $\mathrm{Br}$ & bromo & 80 \\
47 & $\mathrm{Ag}$ & plata & 108 \\
53 & $\mathrm{I}$  & yodo & 127 \\
56 & $\mathrm{Ba}$ & bario & 137 \\
79 & $\mathrm{Au}$ & oro & 197 \\
82 & $\mathrm{Pb}$ & plomo & 207 \\
\hline
\end{tabular}
\end{center}

\subsubsection*{El átomo}

\begin{description}
    \item[$Z$ (número atómico).]
    Número de protones; está en la tabla.

    \item[$A$ (número másico).]
    Protones $+$ neutrones; se escribe arriba a la izquierda:
    ${}^{23}\mathrm{Na}$.

    \item[Neutrones.]
    $A-Z$.

    \item[Electrones de un átomo neutro.]
    $Z$.

    \item[Electrones de un ion.]
    $Z-\mbox{carga}$.
    Para $\mathrm{Na}^{+}$: $11-1=10$;
    para $\mathrm{O}^{2-}$: $8+2=10$.

    \item[Masa atómica media.]
    \[
    \frac{m_1p_1+m_2p_2+\cdots}{100},
    \]
    donde $m_i$ es la masa de cada isótopo y $p_i$ su porcentaje.

    \item[Masa molecular.]
    Suma de las masas atómicas de todos los átomos de la fórmula.
\end{description}

\textbf{Contar átomos y ajustar.}
El subíndice cuenta los átomos del símbolo que tiene delante;
un subíndice detrás de un paréntesis multiplica todo lo de dentro;
el coeficiente multiplica toda la fórmula.
En $\mathrm{Ca(OH)_2}$ hay $1$ Ca, $2$ O y $2$ H;
en $3\,\mathrm{Ca(OH)_2}$ hay $3$ Ca, $6$ O y $6$ H.
Ajustar es poner coeficientes delante de las fórmulas para que cada
elemento tenga el mismo número de átomos a los dos lados de la flecha.
Solo se cambian coeficientes, nunca subíndices, y se dan los enteros
más pequeños posibles. Si un coeficiente es $1$ no se escribe.

\subsubsection*{El mol y los gases}

\begin{description}
    \item[Masa molar $M$.]
    La masa molecular en $\mathrm{g/mol}$:
    $\mathrm{H_2O}$: $18\,\mathrm{u}\rightarrow18\,\mathrm{g/mol}$.

    \item[Moles a partir de la masa.]
    \[
    n=\frac{m}{M}.
    \]

    \item[Partículas a partir de los moles.]
    \[
    N=nN_{\mathrm{A}},
    \qquad N_{\mathrm{A}}=6.022\cdot10^{23}.
    \]

    \item[Volumen de un gas en condiciones normales.]
    A $0\,{}^\circ\mathrm{C}$ y $1\,\mathrm{atm}$:
    \[
    V=n\cdot22.4\,\mathrm{L/mol}.
    \]

    \item[Reacción ajustada.]
    Los coeficientes dan la proporción en moles;
    entre gases a la misma presión y temperatura, también en litros.
\end{description}

\subsubsection*{Disoluciones y rendimiento (4.º)}

\begin{description}
    \item[Concentración en masa.]
    \[
    c=\frac{m(\mbox{soluto})}{V(\mbox{disolución})},
    \qquad \mbox{en }\mathrm{g/L}.
    \]

    \item[Molaridad.]
    \[
    M=\frac{n(\mbox{soluto})}{V(\mbox{disolución en L})},
    \qquad \mbox{en }\mathrm{mol/L}.
    \]

    \item[Porcentaje en masa.]
    \[
    \%=\frac{m(\mbox{soluto})}{m(\mbox{disolución})}\cdot100.
    \]

    \item[Rendimiento.]
    \[
    \%=\frac{\mbox{cantidad obtenida}}{\mbox{cantidad teórica}}\cdot100.
    \]

    \item[Pureza.]
    \[
    \%=\frac{\mbox{masa de sustancia pura}}{\mbox{masa de la muestra}}\cdot100.
    \]

    \item[Reactivo limitante.]
    El que se acaba primero: divide los moles de cada reactivo entre
    su coeficiente; el menor es el limitante.
\end{description}

\subsubsection*{Configuración electrónica (4.º)}

Los electrones se colocan en este orden:
\[
1\mathrm{s}\quad2\mathrm{s}\quad2\mathrm{p}\quad
3\mathrm{s}\quad3\mathrm{p}\quad4\mathrm{s}\quad
3\mathrm{d}\quad4\mathrm{p}\quad5\mathrm{s}.
\]
Caben $2$ en s, $6$ en p y $10$ en d.
Ejemplo: Na ($Z=11$):
\[
1\mathrm{s}^2\ 2\mathrm{s}^2\ 2\mathrm{p}^6\ 3\mathrm{s}^1.
\]

\begin{center}
\begin{tabular}{lrrl}
\hline
Final & Grupo & $e^{-}$ de valencia & Ion que forma \\
\hline
$n\mathrm{s}^1$ & 1 & 1 & Pierde 1: $+1$ \\
$n\mathrm{s}^2$ & 2 & 2 & Pierde 2: $+2$ \\
$n\mathrm{s}^2\,n\mathrm{p}^1$ & 13 & 3 & Pierde 3: $+3$ \\
$n\mathrm{s}^2\,n\mathrm{p}^2$ & 14 & 4 & No se pide \\
$n\mathrm{s}^2\,n\mathrm{p}^3$ & 15 & 5 & Gana 3: $-3$ \\
$n\mathrm{s}^2\,n\mathrm{p}^4$ & 16 & 6 & Gana 2: $-2$ \\
$n\mathrm{s}^2\,n\mathrm{p}^5$ & 17 & 7 & Gana 1: $-1$ \\
$n\mathrm{s}^2\,n\mathrm{p}^6$ & 18 & 8 & Ninguno \\
\hline
\end{tabular}
\end{center}

El periodo es el número $n$ más alto que aparece en la configuración.
En un compuesto iónico las cargas suman cero:
$\mathrm{Al}^{3+}$ y $\mathrm{O}^{2-}$ dan $\mathrm{Al_2O_3}$,
porque $2\cdot(+3)+3\cdot(-2)=0$.

\subsection{3.º ESO: ejercicios 1--10}

\subsubsection*{Ejercicio 1. Símbolos y nombres}

En a)--f) escribe el símbolo. En g)--l) escribe el nombre.

\begin{enumerate}
    \item[a)] Sodio.
    \item[b)] Potasio.
    \item[c)] Hierro.
    \item[d)] Cobre.
    \item[e)] Plata.
    \item[f)] Plomo.
    \item[g)] $\mathrm{Au}$.
    \item[h)] $\mathrm{Zn}$.
    \item[i)] $\mathrm{Mn}$.
    \item[j)] $\mathrm{S}$.
    \item[k)] $\mathrm{P}$.
    \item[l)] $\mathrm{Ba}$.
\end{enumerate}

\subsubsection*{Ejercicio 2. Contar átomos en una fórmula}

Escribe cuántos átomos hay de cada elemento y cuántos átomos hay en total.

\begin{enumerate}
    \item[a)] $\mathrm{H_2O}$.
    \item[b)] $\mathrm{CO_2}$.
    \item[c)] $\mathrm{H_2SO_4}$.
    \item[d)] $\mathrm{Ca(OH)_2}$.
    \item[e)] $\mathrm{Al_2(SO_4)_3}$.
    \item[f)] $\mathrm{C_6H_{12}O_6}$.
    \item[g)] $\mathrm{NH_4NO_3}$.
    \item[h)] $\mathrm{Fe_2O_3}$.
    \item[i)] $\mathrm{Mg_3(PO_4)_2}$.
    \item[j)] $\mathrm{(NH_4)_2CO_3}$.
    \item[k)] $\mathrm{C_2H_5OH}$.
    \item[l)] $\mathrm{Ca_3(PO_4)_2}$.
\end{enumerate}

\subsubsection*{Ejercicio 3. Contar átomos con coeficiente}

El coeficiente multiplica toda la fórmula.
Escribe cuántos átomos hay de cada elemento y en total.

\begin{enumerate}
    \item[a)] $3\,\mathrm{H_2O}$.
    \item[b)] $2\,\mathrm{CO_2}$.
    \item[c)] $5\,\mathrm{O_2}$.
    \item[d)] $4\,\mathrm{NH_3}$.
    \item[e)] $2\,\mathrm{Ca(OH)_2}$.
    \item[f)] $3\,\mathrm{Al_2(SO_4)_3}$.
    \item[g)] $6\,\mathrm{C_6H_{12}O_6}$.
    \item[h)] $10\,\mathrm{HNO_3}$.
\end{enumerate}

\subsubsection*{Ejercicio 4. Partículas de un átomo neutro}

Busca $Z$ en la tabla y escribe el número de protones, neutrones y electrones.

\begin{enumerate}
    \item[a)] ${}^{12}\mathrm{C}$.
    \item[b)] ${}^{14}\mathrm{N}$.
    \item[c)] ${}^{16}\mathrm{O}$.
    \item[d)] ${}^{23}\mathrm{Na}$.
    \item[e)] ${}^{35}\mathrm{Cl}$.
    \item[f)] ${}^{37}\mathrm{Cl}$.
    \item[g)] ${}^{56}\mathrm{Fe}$.
    \item[h)] ${}^{40}\mathrm{Ca}$.
    \item[i)] ${}^{27}\mathrm{Al}$.
    \item[j)] ${}^{197}\mathrm{Au}$.
\end{enumerate}

\subsubsection*{Ejercicio 5. Partículas de un ion}

Escribe el número de protones, neutrones y electrones.

\begin{enumerate}
    \item[a)] ${}^{23}\mathrm{Na}^{+}$.
    \item[b)] ${}^{35}\mathrm{Cl}^{-}$.
    \item[c)] ${}^{24}\mathrm{Mg}^{2+}$.
    \item[d)] ${}^{16}\mathrm{O}^{2-}$.
    \item[e)] ${}^{27}\mathrm{Al}^{3+}$.
    \item[f)] ${}^{32}\mathrm{S}^{2-}$.
    \item[g)] ${}^{56}\mathrm{Fe}^{3+}$.
    \item[h)] ${}^{40}\mathrm{Ca}^{2+}$.
    \item[i)] ${}^{19}\mathrm{F}^{-}$.
    \item[j)] ${}^{39}\mathrm{K}^{+}$.
\end{enumerate}

\subsubsection*{Ejercicio 6. De las partículas al símbolo}

Con el número de protones busca el elemento en la tabla.
Escribe el símbolo con $A$ arriba a la izquierda y la carga arriba a la derecha.

\begin{enumerate}
    \item[a)] $p=11$, $n=12$, $e=10$.
    \item[b)] $p=17$, $n=18$, $e=18$.
    \item[c)] $p=8$, $n=8$, $e=8$.
    \item[d)] $p=26$, $n=30$, $e=23$.
    \item[e)] $p=16$, $n=16$, $e=18$.
    \item[f)] $p=20$, $n=20$, $e=18$.
    \item[g)] $p=13$, $n=14$, $e=10$.
    \item[h)] $p=35$, $n=45$, $e=36$.
\end{enumerate}

\subsubsection*{Ejercicio 7. Isótopos}

En a)--f) escribe protones, neutrones y electrones.
En g)--k) calcula la masa atómica media,
tomando la masa de cada isótopo igual a su $A$.

\begin{enumerate}
    \item[a)] ${}^{1}\mathrm{H}$.
    \item[b)] ${}^{2}\mathrm{H}$.
    \item[c)] ${}^{3}\mathrm{H}$.
    \item[d)] ${}^{12}\mathrm{C}$.
    \item[e)] ${}^{13}\mathrm{C}$.
    \item[f)] ${}^{14}\mathrm{C}$.
    \item[g)] Cloro:
    ${}^{35}\mathrm{Cl}$ ($75.77\,\%$) y
    ${}^{37}\mathrm{Cl}$ ($24.23\,\%$).
    \item[h)] Boro:
    ${}^{10}\mathrm{B}$ ($19.9\,\%$) y
    ${}^{11}\mathrm{B}$ ($80.1\,\%$).
    \item[i)] Cobre:
    ${}^{63}\mathrm{Cu}$ ($69.2\,\%$) y
    ${}^{65}\mathrm{Cu}$ ($30.8\,\%$).
    \item[j)] Magnesio:
    ${}^{24}\mathrm{Mg}$ ($79\,\%$),
    ${}^{25}\mathrm{Mg}$ ($10\,\%$) y
    ${}^{26}\mathrm{Mg}$ ($11\,\%$).
    \item[k)] Litio:
    ${}^{6}\mathrm{Li}$ ($7.5\,\%$) y
    ${}^{7}\mathrm{Li}$ ($92.5\,\%$).
\end{enumerate}

\subsubsection*{Ejercicio 8. Masa molecular}

Calcula la masa molecular en $\mathrm{u}$ con las masas de la tabla.

\begin{enumerate}
    \item[a)] $\mathrm{H_2O}$.
    \item[b)] $\mathrm{CO_2}$.
    \item[c)] $\mathrm{NaCl}$.
    \item[d)] $\mathrm{H_2SO_4}$.
    \item[e)] $\mathrm{CaCO_3}$.
    \item[f)] $\mathrm{NH_3}$.
    \item[g)] $\mathrm{C_6H_{12}O_6}$.
    \item[h)] $\mathrm{Ca(OH)_2}$.
    \item[i)] $\mathrm{Al_2(SO_4)_3}$.
    \item[j)] $\mathrm{HNO_3}$.
    \item[k)] $\mathrm{Fe_2O_3}$.
    \item[l)] $\mathrm{CuSO_4}$.
    \item[m)] $\mathrm{KMnO_4}$.
    \item[n)] $\mathrm{C_2H_5OH}$.
    \item[o)] $\mathrm{Mg_3(PO_4)_2}$.
\end{enumerate}

\subsubsection*{Ejercicio 9. ¿Está ajustada?}

Cuenta los átomos de cada elemento a la izquierda y a la derecha de la flecha.
Responde sí o no. En las soluciones, «H $4|4$» significa
$4$ átomos de H a la izquierda y $4$ a la derecha.

\begin{enumerate}
    \item[a)] $2\,\mathrm{H_2}+\mathrm{O_2}\rightarrow2\,\mathrm{H_2O}$.
    \item[b)] $\mathrm{H_2}+\mathrm{O_2}\rightarrow\mathrm{H_2O}$.
    \item[c)] $\mathrm{CH_4}+2\,\mathrm{O_2}\rightarrow
    \mathrm{CO_2}+2\,\mathrm{H_2O}$.
    \item[d)] $\mathrm{N_2}+\mathrm{H_2}\rightarrow2\,\mathrm{NH_3}$.
    \item[e)] $2\,\mathrm{Na}+\mathrm{Cl_2}\rightarrow2\,\mathrm{NaCl}$.
    \item[f)] $\mathrm{Fe_2O_3}+3\,\mathrm{CO}\rightarrow
    2\,\mathrm{Fe}+2\,\mathrm{CO_2}$.
    \item[g)] $\mathrm{C_3H_8}+5\,\mathrm{O_2}\rightarrow
    3\,\mathrm{CO_2}+4\,\mathrm{H_2O}$.
    \item[h)] $2\,\mathrm{Al}+3\,\mathrm{O_2}\rightarrow\mathrm{Al_2O_3}$.
\end{enumerate}

\subsubsection*{Ejercicio 10. Ajustar reacciones (nivel 1)}

Ajusta cada reacción.

\begin{enumerate}
    \item[a)] $\mathrm{H_2}+\mathrm{O_2}\rightarrow\mathrm{H_2O}$.
    \item[b)] $\mathrm{N_2}+\mathrm{H_2}\rightarrow\mathrm{NH_3}$.
    \item[c)] $\mathrm{Na}+\mathrm{Cl_2}\rightarrow\mathrm{NaCl}$.
    \item[d)] $\mathrm{Mg}+\mathrm{O_2}\rightarrow\mathrm{MgO}$.
    \item[e)] $\mathrm{H_2}+\mathrm{Cl_2}\rightarrow\mathrm{HCl}$.
    \item[f)] $\mathrm{Fe}+\mathrm{S}\rightarrow\mathrm{FeS}$.
    \item[g)] $\mathrm{K}+\mathrm{O_2}\rightarrow\mathrm{K_2O}$.
    \item[h)] $\mathrm{Al}+\mathrm{Cl_2}\rightarrow\mathrm{AlCl_3}$.
    \item[i)] $\mathrm{CaCO_3}\rightarrow\mathrm{CaO}+\mathrm{CO_2}$.
    \item[j)] $\mathrm{H_2O_2}\rightarrow\mathrm{H_2O}+\mathrm{O_2}$.
    \item[k)] $\mathrm{Zn}+\mathrm{HCl}\rightarrow\mathrm{ZnCl_2}+\mathrm{H_2}$.
    \item[l)] $\mathrm{CH_4}+\mathrm{O_2}\rightarrow\mathrm{CO_2}+\mathrm{H_2O}$.
\end{enumerate}

\subsection{3.º ESO: ejercicios 11--20}

\subsubsection*{Ejercicio 11. Ajustar reacciones (nivel 2)}

Ajusta cada reacción.

\begin{enumerate}
    \item[a)] $\mathrm{Fe}+\mathrm{O_2}\rightarrow\mathrm{Fe_2O_3}$.
    \item[b)] $\mathrm{Al}+\mathrm{O_2}\rightarrow\mathrm{Al_2O_3}$.
    \item[c)] $\mathrm{C_2H_6}+\mathrm{O_2}\rightarrow\mathrm{CO_2}+\mathrm{H_2O}$.
    \item[d)] $\mathrm{C_3H_8}+\mathrm{O_2}\rightarrow\mathrm{CO_2}+\mathrm{H_2O}$.
    \item[e)] $\mathrm{Na}+\mathrm{H_2O}\rightarrow\mathrm{NaOH}+\mathrm{H_2}$.
    \item[f)] $\mathrm{Mg}+\mathrm{HCl}\rightarrow\mathrm{MgCl_2}+\mathrm{H_2}$.
    \item[g)] $\mathrm{Fe_2O_3}+\mathrm{CO}\rightarrow\mathrm{Fe}+\mathrm{CO_2}$.
    \item[h)] $\mathrm{KClO_3}\rightarrow\mathrm{KCl}+\mathrm{O_2}$.
    \item[i)] $\mathrm{N_2}+\mathrm{O_2}\rightarrow\mathrm{NO_2}$.
    \item[j)] $\mathrm{P}+\mathrm{O_2}\rightarrow\mathrm{P_2O_5}$.
    \item[k)] $\mathrm{CaCO_3}+\mathrm{HCl}\rightarrow
    \mathrm{CaCl_2}+\mathrm{CO_2}+\mathrm{H_2O}$.
    \item[l)] $\mathrm{Al}+\mathrm{HCl}\rightarrow\mathrm{AlCl_3}+\mathrm{H_2}$.
\end{enumerate}

\subsubsection*{Ejercicio 12. Ajustar reacciones (nivel 3)}

Ajusta cada reacción. Truco para las combustiones:
ajusta primero C, después H y deja el O para el final;
si te sale un coeficiente con media unidad, multiplica todo por $2$.

\begin{enumerate}
    \item[a)] $\mathrm{C_4H_{10}}+\mathrm{O_2}\rightarrow
    \mathrm{CO_2}+\mathrm{H_2O}$.
    \item[b)] $\mathrm{C_2H_5OH}+\mathrm{O_2}\rightarrow
    \mathrm{CO_2}+\mathrm{H_2O}$.
    \item[c)] $\mathrm{C_6H_{12}O_6}+\mathrm{O_2}\rightarrow
    \mathrm{CO_2}+\mathrm{H_2O}$.
    \item[d)] $\mathrm{Fe}+\mathrm{H_2O}\rightarrow\mathrm{Fe_3O_4}+\mathrm{H_2}$.
    \item[e)] $\mathrm{NH_3}+\mathrm{O_2}\rightarrow\mathrm{NO}+\mathrm{H_2O}$.
    \item[f)] $\mathrm{Al}+\mathrm{H_2SO_4}\rightarrow
    \mathrm{Al_2(SO_4)_3}+\mathrm{H_2}$.
    \item[g)] $\mathrm{Ca(OH)_2}+\mathrm{H_3PO_4}\rightarrow
    \mathrm{Ca_3(PO_4)_2}+\mathrm{H_2O}$.
    \item[h)] $\mathrm{FeS_2}+\mathrm{O_2}\rightarrow
    \mathrm{Fe_2O_3}+\mathrm{SO_2}$.
    \item[i)] $\mathrm{C_8H_{18}}+\mathrm{O_2}\rightarrow
    \mathrm{CO_2}+\mathrm{H_2O}$.
    \item[j)] $\mathrm{Pb(NO_3)_2}\rightarrow
    \mathrm{PbO}+\mathrm{NO_2}+\mathrm{O_2}$.
\end{enumerate}

\subsubsection*{Ejercicio 13. Conservación de la masa}

La masa total de los reactivos es igual a la masa total de los productos.
Calcula la masa que falta.

\begin{enumerate}
    \item[a)] $4\,\mathrm{g}$ de $\mathrm{H_2}$ reaccionan con
    $32\,\mathrm{g}$ de $\mathrm{O_2}$ y se forma agua.
    Masa de agua.

    \item[b)] $12\,\mathrm{g}$ de C reaccionan con
    $32\,\mathrm{g}$ de $\mathrm{O_2}$ y se forma $\mathrm{CO_2}$.
    Masa de $\mathrm{CO_2}$.

    \item[c)] $100\,\mathrm{g}$ de $\mathrm{CaCO_3}$ se descomponen en
    $56\,\mathrm{g}$ de $\mathrm{CaO}$ y $\mathrm{CO_2}$.
    Masa de $\mathrm{CO_2}$.

    \item[d)] $56\,\mathrm{g}$ de Fe reaccionan con azufre
    y se forman $88\,\mathrm{g}$ de $\mathrm{FeS}$.
    Masa de azufre.

    \item[e)] $2.4\,\mathrm{g}$ de Mg reaccionan con oxígeno
    y se forman $4.0\,\mathrm{g}$ de $\mathrm{MgO}$.
    Masa de oxígeno.

    \item[f)] Zn reacciona con $7.3\,\mathrm{g}$ de $\mathrm{HCl}$
    y se forman $13.6\,\mathrm{g}$ de $\mathrm{ZnCl_2}$
    y $0.2\,\mathrm{g}$ de $\mathrm{H_2}$.
    Masa de Zn.

    \item[g)] $16\,\mathrm{g}$ de $\mathrm{CH_4}$ reaccionan con
    $64\,\mathrm{g}$ de $\mathrm{O_2}$ y se forman
    $44\,\mathrm{g}$ de $\mathrm{CO_2}$ y agua.
    Masa de agua.

    \item[h)] $23\,\mathrm{g}$ de Na reaccionan con
    $18\,\mathrm{g}$ de $\mathrm{H_2O}$ y se forman
    $40\,\mathrm{g}$ de $\mathrm{NaOH}$ e hidrógeno.
    Masa de hidrógeno.
\end{enumerate}

\subsubsection*{Ejercicio 14. Proporciones fijas en masa}

Un compuesto tiene siempre la misma proporción en masa entre sus elementos.
Usa la proporción que se da.

\begin{enumerate}
    \item[a)] Agua ($\mathrm{H}:\mathrm{O}=1:8$).
    ¿Cuántos gramos de O se combinan con $3\,\mathrm{g}$ de H?
    \item[b)] Agua.
    ¿Cuántos gramos de H se combinan con $40\,\mathrm{g}$ de O?
    \item[c)] Agua.
    ¿Cuántos gramos de agua se forman con $2\,\mathrm{g}$ de H?
    \item[d)] $\mathrm{CO_2}$ ($\mathrm{C}:\mathrm{O}=3:8$).
    ¿Cuántos gramos de O se combinan con $6\,\mathrm{g}$ de C?
    \item[e)] $\mathrm{CO_2}$.
    ¿Cuántos gramos de C hay en $110\,\mathrm{g}$ de $\mathrm{CO_2}$?
    \item[f)] $\mathrm{FeS}$ ($\mathrm{Fe}:\mathrm{S}=7:4$).
    ¿Cuántos gramos de S se combinan con $28\,\mathrm{g}$ de Fe?
    \item[g)] $\mathrm{FeS}$.
    Se mezclan $35\,\mathrm{g}$ de Fe y $35\,\mathrm{g}$ de S.
    ¿Cuál sobra, cuántos gramos sobran y cuántos gramos de $\mathrm{FeS}$
    se forman?
\end{enumerate}

\subsubsection*{Ejercicio 15. Moles y masa}

Calcula la masa molar y después los moles ($n=m/M$)
o la masa ($m=nM$).

\begin{enumerate}
    \item[a)] $36\,\mathrm{g}$ de $\mathrm{H_2O}\rightarrow\mathrm{mol}$.
    \item[b)] $88\,\mathrm{g}$ de $\mathrm{CO_2}\rightarrow\mathrm{mol}$.
    \item[c)] $0.5\,\mathrm{mol}$ de $\mathrm{NaCl}\rightarrow\mathrm{g}$.
    \item[d)] $3\,\mathrm{mol}$ de $\mathrm{O_2}\rightarrow\mathrm{g}$.
    \item[e)] $4.9\,\mathrm{g}$ de $\mathrm{H_2SO_4}\rightarrow\mathrm{mol}$.
    \item[f)] $0.25\,\mathrm{mol}$ de $\mathrm{CaCO_3}\rightarrow\mathrm{g}$.
    \item[g)] $10\,\mathrm{g}$ de $\mathrm{NaOH}\rightarrow\mathrm{mol}$.
    \item[h)] $1.5\,\mathrm{mol}$ de $\mathrm{C_6H_{12}O_6}\rightarrow\mathrm{g}$.
    \item[i)] $171\,\mathrm{g}$ de $\mathrm{Al_2(SO_4)_3}\rightarrow\mathrm{mol}$.
    \item[j)] $2.8\,\mathrm{g}$ de $\mathrm{Fe}\rightarrow\mathrm{mol}$.
    \item[k)] $0.2\,\mathrm{mol}$ de $\mathrm{KMnO_4}\rightarrow\mathrm{g}$.
    \item[l)] $1\,\mathrm{kg}$ de $\mathrm{CaCO_3}\rightarrow\mathrm{mol}$.
\end{enumerate}

\subsubsection*{Ejercicio 16. Moles y partículas}

Usa $N=nN_{\mathrm{A}}$ con $N_{\mathrm{A}}=6.022\cdot10^{23}$.

\begin{enumerate}
    \item[a)] $2\,\mathrm{mol}$ de $\mathrm{H_2O}\rightarrow$ moléculas.
    \item[b)] $0.5\,\mathrm{mol}$ de Fe $\rightarrow$ átomos.
    \item[c)] $18\,\mathrm{g}$ de $\mathrm{H_2O}\rightarrow$
    moles, moléculas y átomos de H.
    \item[d)] $44\,\mathrm{g}$ de $\mathrm{CO_2}\rightarrow$
    moles, moléculas y átomos de O.
    \item[e)] $1.2044\cdot10^{24}$ moléculas de $\mathrm{O_2}\rightarrow$
    moles y gramos.
    \item[f)] $3.011\cdot10^{23}$ átomos de Na $\rightarrow$
    moles y gramos.
    \item[g)] $1\,\mathrm{mol}$ de $\mathrm{NH_3}\rightarrow$ átomos de H.
    \item[h)] $6.4\,\mathrm{g}$ de $\mathrm{O_2}\rightarrow$
    moles y moléculas.
\end{enumerate}

\subsubsection*{Ejercicio 17. Moles en una reacción ajustada}

Los coeficientes son la proporción en moles.

\begin{enumerate}
    \item[a)] $\mathrm{N_2}+3\,\mathrm{H_2}\rightarrow2\,\mathrm{NH_3}$.
    Con $2\,\mathrm{mol}$ de $\mathrm{N_2}$:
    moles de $\mathrm{H_2}$ que reaccionan y moles de $\mathrm{NH_3}$
    que se forman.

    \item[b)] $\mathrm{N_2}+3\,\mathrm{H_2}\rightarrow2\,\mathrm{NH_3}$.
    Con $9\,\mathrm{mol}$ de $\mathrm{H_2}$:
    moles de $\mathrm{N_2}$ y de $\mathrm{NH_3}$.

    \item[c)] $\mathrm{N_2}+3\,\mathrm{H_2}\rightarrow2\,\mathrm{NH_3}$.
    Para obtener $10\,\mathrm{mol}$ de $\mathrm{NH_3}$:
    moles de $\mathrm{N_2}$ y de $\mathrm{H_2}$.

    \item[d)] $\mathrm{CH_4}+2\,\mathrm{O_2}\rightarrow
    \mathrm{CO_2}+2\,\mathrm{H_2O}$.
    Con $3\,\mathrm{mol}$ de $\mathrm{CH_4}$:
    moles de $\mathrm{O_2}$, $\mathrm{CO_2}$ y $\mathrm{H_2O}$.

    \item[e)] $\mathrm{CH_4}+2\,\mathrm{O_2}\rightarrow
    \mathrm{CO_2}+2\,\mathrm{H_2O}$.
    Con $5\,\mathrm{mol}$ de $\mathrm{O_2}$:
    moles de $\mathrm{CH_4}$ y de $\mathrm{H_2O}$.

    \item[f)] $2\,\mathrm{Al}+3\,\mathrm{Cl_2}\rightarrow2\,\mathrm{AlCl_3}$.
    Con $4\,\mathrm{mol}$ de Al:
    moles de $\mathrm{Cl_2}$ y de $\mathrm{AlCl_3}$.

    \item[g)] $2\,\mathrm{Al}+3\,\mathrm{Cl_2}\rightarrow2\,\mathrm{AlCl_3}$.
    Con $0.6\,\mathrm{mol}$ de $\mathrm{Cl_2}$:
    moles de Al y de $\mathrm{AlCl_3}$.

    \item[h)] $\mathrm{C_3H_8}+5\,\mathrm{O_2}\rightarrow
    3\,\mathrm{CO_2}+4\,\mathrm{H_2O}$.
    Con $2\,\mathrm{mol}$ de $\mathrm{C_3H_8}$:
    moles de $\mathrm{O_2}$, $\mathrm{CO_2}$ y $\mathrm{H_2O}$.
\end{enumerate}

\subsubsection*{Ejercicio 18. Cálculos con masas}

Ajusta la reacción, pasa la masa dada a moles,
usa la proporción de coeficientes y pasa el resultado a gramos.

\begin{enumerate}
    \item[a)] $\mathrm{H_2}+\mathrm{O_2}\rightarrow\mathrm{H_2O}$:
    masa de $\mathrm{H_2O}$ a partir de $8\,\mathrm{g}$ de $\mathrm{H_2}$.

    \item[b)] $\mathrm{CaCO_3}\rightarrow\mathrm{CaO}+\mathrm{CO_2}$:
    masa de $\mathrm{CaO}$ a partir de $50\,\mathrm{g}$ de $\mathrm{CaCO_3}$.

    \item[c)] $\mathrm{CH_4}+\mathrm{O_2}\rightarrow
    \mathrm{CO_2}+\mathrm{H_2O}$:
    masa de $\mathrm{O_2}$ a partir de $32\,\mathrm{g}$ de $\mathrm{CH_4}$.

    \item[d)] $\mathrm{CH_4}+\mathrm{O_2}\rightarrow
    \mathrm{CO_2}+\mathrm{H_2O}$:
    masa de $\mathrm{CO_2}$ a partir de $32\,\mathrm{g}$ de $\mathrm{CH_4}$.

    \item[e)] $\mathrm{Mg}+\mathrm{O_2}\rightarrow\mathrm{MgO}$:
    masa de $\mathrm{MgO}$ a partir de $12\,\mathrm{g}$ de Mg.

    \item[f)] $\mathrm{N_2}+\mathrm{H_2}\rightarrow\mathrm{NH_3}$:
    masa de $\mathrm{NH_3}$ a partir de $56\,\mathrm{g}$ de $\mathrm{N_2}$.

    \item[g)] $\mathrm{Fe_2O_3}+\mathrm{CO}\rightarrow\mathrm{Fe}+\mathrm{CO_2}$:
    masa de Fe a partir de $320\,\mathrm{g}$ de $\mathrm{Fe_2O_3}$.

    \item[h)] $\mathrm{KClO_3}\rightarrow\mathrm{KCl}+\mathrm{O_2}$:
    masa de $\mathrm{O_2}$ a partir de $24.5\,\mathrm{g}$ de $\mathrm{KClO_3}$.

    \item[i)] $\mathrm{Zn}+\mathrm{HCl}\rightarrow\mathrm{ZnCl_2}+\mathrm{H_2}$:
    masa de $\mathrm{H_2}$ a partir de $13.08\,\mathrm{g}$ de Zn.
\end{enumerate}

\subsubsection*{Ejercicio 19. Volumen de gases en condiciones normales}

Todos los gases están a $0\,{}^\circ\mathrm{C}$ y $1\,\mathrm{atm}$:
$1\,\mathrm{mol}$ ocupa $22.4\,\mathrm{L}$.

\begin{enumerate}
    \item[a)] $2\,\mathrm{mol}$ de $\mathrm{O_2}\rightarrow$ litros.
    \item[b)] $11.2\,\mathrm{L}$ de $\mathrm{CO_2}\rightarrow$ moles y gramos.
    \item[c)] $64\,\mathrm{g}$ de $\mathrm{O_2}\rightarrow$ litros.
    \item[d)] $5.6\,\mathrm{L}$ de $\mathrm{N_2}\rightarrow$ gramos.
    \item[e)] $1\,\mathrm{g}$ de $\mathrm{H_2}\rightarrow$ litros.
    \item[f)] $67.2\,\mathrm{L}$ de $\mathrm{NH_3}\rightarrow$ moles y gramos.
    \item[g)] $4.4\,\mathrm{g}$ de $\mathrm{CO_2}\rightarrow$ litros.
    \item[h)] $3.011\cdot10^{23}$ moléculas de $\mathrm{CH_4}\rightarrow$ litros.
\end{enumerate}

\subsubsection*{Ejercicio 20. Volúmenes en las reacciones}

En a)--d) los gases están a la misma presión y temperatura:
los coeficientes dan la proporción en litros.
En e)--h) pasa la masa a moles y usa $22.4\,\mathrm{L/mol}$.

\begin{enumerate}
    \item[a)] $\mathrm{N_2}+3\,\mathrm{H_2}\rightarrow2\,\mathrm{NH_3}$
    (todo gas).
    Con $10\,\mathrm{L}$ de $\mathrm{N_2}$:
    litros de $\mathrm{H_2}$ y de $\mathrm{NH_3}$.

    \item[b)] $2\,\mathrm{H_2}+\mathrm{O_2}\rightarrow2\,\mathrm{H_2O}$
    (todo gas).
    Con $6\,\mathrm{L}$ de $\mathrm{H_2}$:
    litros de $\mathrm{O_2}$ y de $\mathrm{H_2O}$.

    \item[c)] $\mathrm{CH_4}+2\,\mathrm{O_2}\rightarrow
    \mathrm{CO_2}+2\,\mathrm{H_2O}$ (todo gas).
    Con $5\,\mathrm{L}$ de $\mathrm{CH_4}$:
    litros de $\mathrm{O_2}$, $\mathrm{CO_2}$ y $\mathrm{H_2O}$.

    \item[d)] $\mathrm{H_2}+\mathrm{Cl_2}\rightarrow2\,\mathrm{HCl}$
    (todo gas).
    Con $15\,\mathrm{L}$ de $\mathrm{Cl_2}$:
    litros de $\mathrm{H_2}$ y de $\mathrm{HCl}$.

    \item[e)] $2\,\mathrm{CO}+\mathrm{O_2}\rightarrow2\,\mathrm{CO_2}$.
    Con $28\,\mathrm{g}$ de $\mathrm{CO}$:
    litros de $\mathrm{O_2}$ en condiciones normales.

    \item[f)] $\mathrm{CaCO_3}\rightarrow\mathrm{CaO}+\mathrm{CO_2}$.
    Con $50\,\mathrm{g}$ de $\mathrm{CaCO_3}$:
    litros de $\mathrm{CO_2}$ en condiciones normales.

    \item[g)] $\mathrm{Zn}+2\,\mathrm{HCl}\rightarrow
    \mathrm{ZnCl_2}+\mathrm{H_2}$.
    Con $6.54\,\mathrm{g}$ de Zn:
    litros de $\mathrm{H_2}$ en condiciones normales.

    \item[h)] $\mathrm{C_3H_8}+5\,\mathrm{O_2}\rightarrow
    3\,\mathrm{CO_2}+4\,\mathrm{H_2O}$.
    Con $22\,\mathrm{g}$ de $\mathrm{C_3H_8}$:
    litros de $\mathrm{O_2}$ y de $\mathrm{CO_2}$ en condiciones normales.
\end{enumerate}

\subsection{4.º ESO: ejercicios 21--30}

\subsubsection*{Ejercicio 21. Configuración electrónica}

Escribe la configuración electrónica del elemento,
siguiendo el orden de llenado de la tabla.

\begin{enumerate}
    \item[a)] H ($Z=1$).
    \item[b)] C ($Z=6$).
    \item[c)] N ($Z=7$).
    \item[d)] O ($Z=8$).
    \item[e)] F ($Z=9$).
    \item[f)] Ne ($Z=10$).
    \item[g)] Na ($Z=11$).
    \item[h)] Al ($Z=13$).
    \item[i)] S ($Z=16$).
    \item[j)] Cl ($Z=17$).
    \item[k)] K ($Z=19$).
    \item[l)] Ca ($Z=20$).
    \item[m)] Fe ($Z=26$).
    \item[n)] Br ($Z=35$).
\end{enumerate}

\subsubsection*{Ejercicio 22. Grupo y periodo}

Escribe la configuración, y a partir de ella el grupo,
el periodo y los electrones de valencia.

\begin{enumerate}
    \item[a)] $Z=3$.
    \item[b)] $Z=7$.
    \item[c)] $Z=9$.
    \item[d)] $Z=12$.
    \item[e)] $Z=13$.
    \item[f)] $Z=15$.
    \item[g)] $Z=16$.
    \item[h)] $Z=17$.
    \item[i)] $Z=18$.
    \item[j)] $Z=19$.
    \item[k)] $Z=20$.
    \item[l)] $Z=35$.
\end{enumerate}

\subsubsection*{Ejercicio 23. Ion más probable}

Busca $Z$, escribe la configuración, saca el grupo
y con la tabla del principio escribe el ion que forma
y cuántos electrones tiene.
Comprueba que tiene los mismos electrones que un gas noble
(He $2$, Ne $10$, Ar $18$, Kr $36$).

\begin{enumerate}
    \item[a)] Li.
    \item[b)] Mg.
    \item[c)] Al.
    \item[d)] N.
    \item[e)] O.
    \item[f)] F.
    \item[g)] Na.
    \item[h)] S.
    \item[i)] Cl.
    \item[j)] Ca.
    \item[k)] K.
    \item[l)] Br.
\end{enumerate}

\subsubsection*{Ejercicio 24. Fórmulas de compuestos iónicos}

Escribe la fórmula:
el número de cada ion tiene que hacer que las cargas sumen cero.

\begin{enumerate}
    \item[a)] $\mathrm{Na}^{+}$ y $\mathrm{Cl}^{-}$.
    \item[b)] $\mathrm{Mg}^{2+}$ y $\mathrm{Cl}^{-}$.
    \item[c)] $\mathrm{Al}^{3+}$ y $\mathrm{O}^{2-}$.
    \item[d)] $\mathrm{Ca}^{2+}$ y $\mathrm{O}^{2-}$.
    \item[e)] $\mathrm{K}^{+}$ y $\mathrm{S}^{2-}$.
    \item[f)] $\mathrm{Al}^{3+}$ y $\mathrm{Cl}^{-}$.
    \item[g)] $\mathrm{Fe}^{3+}$ y $\mathrm{O}^{2-}$.
    \item[h)] $\mathrm{Mg}^{2+}$ y $\mathrm{N}^{3-}$.
    \item[i)] $\mathrm{Li}^{+}$ y $\mathrm{O}^{2-}$.
    \item[j)] $\mathrm{Ca}^{2+}$ y $\mathrm{F}^{-}$.
    \item[k)] $\mathrm{Al}^{3+}$ y $\mathrm{S}^{2-}$.
    \item[l)] $\mathrm{Ba}^{2+}$ y $\mathrm{Br}^{-}$.
\end{enumerate}

\subsubsection*{Ejercicio 25. Ajustar reacciones (nivel 4)}

Ajusta cada reacción.
Los dos últimos son largos: escribe una tabla con cada elemento y ve probando.

\begin{enumerate}
    \item[a)] $\mathrm{C_5H_{12}}+\mathrm{O_2}\rightarrow
    \mathrm{CO_2}+\mathrm{H_2O}$.
    \item[b)] $\mathrm{C_3H_8O_3}+\mathrm{O_2}\rightarrow
    \mathrm{CO_2}+\mathrm{H_2O}$.
    \item[c)] $\mathrm{C_{12}H_{22}O_{11}}+\mathrm{O_2}\rightarrow
    \mathrm{CO_2}+\mathrm{H_2O}$.
    \item[d)] $\mathrm{Fe_2O_3}+\mathrm{C}\rightarrow\mathrm{Fe}+\mathrm{CO_2}$.
    \item[e)] $\mathrm{Al}+\mathrm{Fe_2O_3}\rightarrow
    \mathrm{Al_2O_3}+\mathrm{Fe}$.
    \item[f)] $\mathrm{NH_4NO_3}\rightarrow\mathrm{N_2O}+\mathrm{H_2O}$.
    \item[g)] $\mathrm{MnO_2}+\mathrm{HCl}\rightarrow
    \mathrm{MnCl_2}+\mathrm{Cl_2}+\mathrm{H_2O}$.
    \item[h)] $\mathrm{Cu}+\mathrm{HNO_3}\rightarrow
    \mathrm{Cu(NO_3)_2}+\mathrm{NO}+\mathrm{H_2O}$.
    \item[i)] $\mathrm{KMnO_4}+\mathrm{HCl}\rightarrow
    \mathrm{KCl}+\mathrm{MnCl_2}+\mathrm{Cl_2}+\mathrm{H_2O}$.
    \item[j)] $\mathrm{Ca_3(PO_4)_2}+\mathrm{SiO_2}+\mathrm{C}\rightarrow
    \mathrm{CaSiO_3}+\mathrm{P_4}+\mathrm{CO}$.
\end{enumerate}

\subsubsection*{Ejercicio 26. Cálculos completos}

Masas, moles, volúmenes en condiciones normales y partículas mezclados.
Las reacciones ya están ajustadas.

\begin{enumerate}
    \item[a)] $2\,\mathrm{H_2O_2}\rightarrow2\,\mathrm{H_2O}+\mathrm{O_2}$.
    Con $68\,\mathrm{g}$ de $\mathrm{H_2O_2}$:
    litros de $\mathrm{O_2}$.

    \item[b)] $\mathrm{Zn}+2\,\mathrm{HCl}\rightarrow
    \mathrm{ZnCl_2}+\mathrm{H_2}$.
    Con $32.7\,\mathrm{g}$ de Zn:
    litros de $\mathrm{H_2}$.

    \item[c)] $\mathrm{C_3H_8}+5\,\mathrm{O_2}\rightarrow
    3\,\mathrm{CO_2}+4\,\mathrm{H_2O}$.
    Para obtener $67.2\,\mathrm{L}$ de $\mathrm{CO_2}$:
    gramos de $\mathrm{C_3H_8}$.

    \item[d)] $2\,\mathrm{Na}+2\,\mathrm{H_2O}\rightarrow
    2\,\mathrm{NaOH}+\mathrm{H_2}$.
    Con $4.6\,\mathrm{g}$ de Na:
    gramos de $\mathrm{NaOH}$ y litros de $\mathrm{H_2}$.

    \item[e)] $\mathrm{CaCO_3}+2\,\mathrm{HCl}\rightarrow
    \mathrm{CaCl_2}+\mathrm{CO_2}+\mathrm{H_2O}$.
    Con $25\,\mathrm{g}$ de $\mathrm{CaCO_3}$:
    gramos de $\mathrm{CaCl_2}$ y litros de $\mathrm{CO_2}$.

    \item[f)] $4\,\mathrm{Fe}+3\,\mathrm{O_2}\rightarrow2\,\mathrm{Fe_2O_3}$.
    Con $11.2\,\mathrm{g}$ de Fe:
    gramos de $\mathrm{Fe_2O_3}$ y litros de $\mathrm{O_2}$.

    \item[g)] $\mathrm{N_2}+3\,\mathrm{H_2}\rightarrow2\,\mathrm{NH_3}$.
    Con $44.8\,\mathrm{L}$ de $\mathrm{H_2}$:
    gramos de $\mathrm{NH_3}$.

    \item[h)] $\mathrm{C_6H_{12}O_6}+6\,\mathrm{O_2}\rightarrow
    6\,\mathrm{CO_2}+6\,\mathrm{H_2O}$.
    Con $90\,\mathrm{g}$ de glucosa:
    moléculas de $\mathrm{CO_2}$.
\end{enumerate}

\subsubsection*{Ejercicio 27. Reactivo limitante}

Ajusta, pasa cada reactivo a moles, divide entre su coeficiente:
el menor es el limitante.
Calcula la masa del producto indicado y cuánto sobra del otro reactivo.

\begin{enumerate}
    \item[a)] $\mathrm{H_2}+\mathrm{O_2}\rightarrow\mathrm{H_2O}$:
    $4\,\mathrm{g}$ de $\mathrm{H_2}$ y $40\,\mathrm{g}$ de $\mathrm{O_2}$;
    producto: $\mathrm{H_2O}$.

    \item[b)] $\mathrm{N_2}+\mathrm{H_2}\rightarrow\mathrm{NH_3}$:
    $28\,\mathrm{g}$ de $\mathrm{N_2}$ y $9\,\mathrm{g}$ de $\mathrm{H_2}$;
    producto: $\mathrm{NH_3}$.

    \item[c)] $\mathrm{Zn}+\mathrm{HCl}\rightarrow\mathrm{ZnCl_2}+\mathrm{H_2}$:
    $13.08\,\mathrm{g}$ de Zn y $18.25\,\mathrm{g}$ de $\mathrm{HCl}$;
    producto: $\mathrm{H_2}$.

    \item[d)] $\mathrm{Al}+\mathrm{Cl_2}\rightarrow\mathrm{AlCl_3}$:
    $5.4\,\mathrm{g}$ de Al y $21.3\,\mathrm{g}$ de $\mathrm{Cl_2}$;
    producto: $\mathrm{AlCl_3}$.

    \item[e)] $\mathrm{CH_4}+\mathrm{O_2}\rightarrow
    \mathrm{CO_2}+\mathrm{H_2O}$:
    $8\,\mathrm{g}$ de $\mathrm{CH_4}$ y $48\,\mathrm{g}$ de $\mathrm{O_2}$;
    producto: $\mathrm{CO_2}$.

    \item[f)] $\mathrm{Fe}+\mathrm{S}\rightarrow\mathrm{FeS}$:
    $28\,\mathrm{g}$ de Fe y $20\,\mathrm{g}$ de S;
    producto: $\mathrm{FeS}$.
\end{enumerate}

\subsubsection*{Ejercicio 28. Rendimiento y pureza}

Usa las fórmulas de rendimiento y pureza de la tabla.

\begin{enumerate}
    \item[a)] $\mathrm{CaCO_3}\rightarrow\mathrm{CaO}+\mathrm{CO_2}$.
    De $100\,\mathrm{g}$ de $\mathrm{CaCO_3}$ se obtienen
    $50.4\,\mathrm{g}$ de $\mathrm{CaO}$. Rendimiento.

    \item[b)] $\mathrm{N_2}+3\,\mathrm{H_2}\rightarrow2\,\mathrm{NH_3}$.
    De $28\,\mathrm{g}$ de $\mathrm{N_2}$ se obtienen
    $25.5\,\mathrm{g}$ de $\mathrm{NH_3}$. Rendimiento.

    \item[c)] $2\,\mathrm{H_2}+\mathrm{O_2}\rightarrow2\,\mathrm{H_2O}$
    con rendimiento del $80\,\%$.
    Gramos de $\mathrm{H_2O}$ que se obtienen de
    $8\,\mathrm{g}$ de $\mathrm{H_2}$.

    \item[d)] $\mathrm{CaCO_3}\rightarrow\mathrm{CaO}+\mathrm{CO_2}$
    con rendimiento del $60\,\%$.
    Litros de $\mathrm{CO_2}$ en condiciones normales que se obtienen
    de $50\,\mathrm{g}$ de $\mathrm{CaCO_3}$.

    \item[e)] $\mathrm{CaCO_3}\rightarrow\mathrm{CaO}+\mathrm{CO_2}$.
    $200\,\mathrm{g}$ de una caliza con $75\,\%$ de $\mathrm{CaCO_3}$.
    Gramos de $\mathrm{CO_2}$.

    \item[f)] $\mathrm{Zn}+2\,\mathrm{HCl}\rightarrow
    \mathrm{ZnCl_2}+\mathrm{H_2}$.
    $20\,\mathrm{g}$ de una muestra con $65.4\,\%$ de Zn.
    Litros de $\mathrm{H_2}$ en condiciones normales.

    \item[g)] $\mathrm{Fe_2O_3}+3\,\mathrm{CO}\rightarrow
    2\,\mathrm{Fe}+3\,\mathrm{CO_2}$.
    $80\,\mathrm{g}$ de mineral dan $42\,\mathrm{g}$ de Fe.
    Pureza del mineral en $\mathrm{Fe_2O_3}$.

    \item[h)] $\mathrm{CaCO_3}\rightarrow\mathrm{CaO}+\mathrm{CO_2}$.
    $100\,\mathrm{g}$ de mineral del $90\,\%$ de pureza,
    rendimiento del $80\,\%$.
    Gramos de $\mathrm{CaO}$.
\end{enumerate}

\subsubsection*{Ejercicio 29. Disoluciones}

Pasa siempre los volúmenes a litros.

\begin{enumerate}
    \item[a)] $20\,\mathrm{g}$ de $\mathrm{NaCl}$ en
    $500\,\mathrm{mL}$ de disolución:
    concentración en $\mathrm{g/L}$.

    \item[b)] $0.5\,\mathrm{mol}$ de soluto en
    $250\,\mathrm{mL}$ de disolución: molaridad.

    \item[c)] $4\,\mathrm{g}$ de $\mathrm{NaOH}$ en
    $200\,\mathrm{mL}$ de disolución: molaridad.

    \item[d)] $9.8\,\mathrm{g}$ de $\mathrm{H_2SO_4}$ en
    $500\,\mathrm{mL}$ de disolución: molaridad.

    \item[e)] Gramos de $\mathrm{NaCl}$ para preparar
    $250\,\mathrm{mL}$ de disolución $0.2\,\mathrm{M}$.

    \item[f)] $15\,\mathrm{g}$ de azúcar disueltos en
    $135\,\mathrm{g}$ de agua: porcentaje en masa.

    \item[g)] Gramos de soluto en $300\,\mathrm{g}$ de disolución
    al $5\,\%$ en masa.

    \item[h)] $\mathrm{mL}$ de disolución $0.5\,\mathrm{M}$
    que contienen $0.1\,\mathrm{mol}$ de soluto.

    \item[i)] $50\,\mathrm{mL}$ de disolución $2\,\mathrm{M}$
    se diluyen con agua hasta $250\,\mathrm{mL}$:
    molaridad final (los moles no cambian).

    \item[j)] Disolución de $\mathrm{NaCl}$ $0.5\,\mathrm{M}$:
    concentración en $\mathrm{g/L}$.

    \item[k)] Gramos de glucosa ($\mathrm{C_6H_{12}O_6}$)
    para preparar $2\,\mathrm{L}$ de disolución $0.1\,\mathrm{M}$.

    \item[l)] Molaridad de una disolución de
    $40\,\mathrm{g/L}$ de $\mathrm{NaOH}$.
\end{enumerate}

\subsubsection*{Ejercicio 30. Reacciones en disolución}

Moles de soluto: $n=MV$, con $V$ en litros.
Los gases, en condiciones normales.

\begin{enumerate}
    \item[a)] $\mathrm{Zn}+2\,\mathrm{HCl}\rightarrow
    \mathrm{ZnCl_2}+\mathrm{H_2}$.
    $\mathrm{mL}$ de $\mathrm{HCl}$ $2\,\mathrm{M}$
    para que reaccionen $6.54\,\mathrm{g}$ de Zn.

    \item[b)] $\mathrm{NaOH}+\mathrm{HCl}\rightarrow
    \mathrm{NaCl}+\mathrm{H_2O}$.
    $\mathrm{mL}$ de $\mathrm{HCl}$ $0.1\,\mathrm{M}$
    para neutralizar $25\,\mathrm{mL}$ de $\mathrm{NaOH}$ $0.2\,\mathrm{M}$.

    \item[c)] $\mathrm{H_2SO_4}+2\,\mathrm{NaOH}\rightarrow
    \mathrm{Na_2SO_4}+2\,\mathrm{H_2O}$.
    $\mathrm{mL}$ de $\mathrm{NaOH}$ $1\,\mathrm{M}$
    para neutralizar $20\,\mathrm{mL}$ de $\mathrm{H_2SO_4}$ $0.5\,\mathrm{M}$.

    \item[d)] $\mathrm{CaCO_3}+2\,\mathrm{HCl}\rightarrow
    \mathrm{CaCl_2}+\mathrm{CO_2}+\mathrm{H_2O}$.
    Con $100\,\mathrm{mL}$ de $\mathrm{HCl}$ $1\,\mathrm{M}$:
    gramos de $\mathrm{CaCO_3}$ que reaccionan y litros de $\mathrm{CO_2}$.

    \item[e)] $\mathrm{Mg}+2\,\mathrm{HCl}\rightarrow
    \mathrm{MgCl_2}+\mathrm{H_2}$.
    $2.4\,\mathrm{g}$ de Mg con $200\,\mathrm{mL}$ de $\mathrm{HCl}$
    $0.5\,\mathrm{M}$:
    limitante, litros de $\mathrm{H_2}$ y gramos que sobran.

    \item[f)] $\mathrm{AgNO_3}+\mathrm{NaCl}\rightarrow
    \mathrm{AgCl}+\mathrm{NaNO_3}$.
    Con $50\,\mathrm{mL}$ de $\mathrm{AgNO_3}$ $0.1\,\mathrm{M}$
    y $\mathrm{NaCl}$ de sobra: gramos de $\mathrm{AgCl}$.
\end{enumerate}

\subsection{Soluciones de 3.º ESO}

\subsubsection*{1. Símbolos y nombres}
\begin{enumerate}
    \item[a)] $\mathrm{Na}$.
    \item[b)] $\mathrm{K}$.
    \item[c)] $\mathrm{Fe}$.
    \item[d)] $\mathrm{Cu}$.
    \item[e)] $\mathrm{Ag}$.
    \item[f)] $\mathrm{Pb}$.
    \item[g)] Oro.
    \item[h)] Cinc.
    \item[i)] Manganeso.
    \item[j)] Azufre.
    \item[k)] Fósforo.
    \item[l)] Bario.
\end{enumerate}

\subsubsection*{2. Contar átomos en una fórmula}
\begin{enumerate}
    \item[a)] H: $2$, O: $1$; total $3$.
    \item[b)] C: $1$, O: $2$; total $3$.
    \item[c)] H: $2$, S: $1$, O: $4$; total $7$.
    \item[d)] Ca: $1$, O: $2$, H: $2$; total $5$.
    \item[e)] Al: $2$, S: $3$, O: $12$; total $17$.
    \item[f)] C: $6$, H: $12$, O: $6$; total $24$.
    \item[g)] N: $2$, H: $4$, O: $3$; total $9$.
    \item[h)] Fe: $2$, O: $3$; total $5$.
    \item[i)] Mg: $3$, P: $2$, O: $8$; total $13$.
    \item[j)] N: $2$, H: $8$, C: $1$, O: $3$; total $14$.
    \item[k)] C: $2$, H: $6$, O: $1$; total $9$.
    \item[l)] Ca: $3$, P: $2$, O: $8$; total $13$.
\end{enumerate}

\subsubsection*{3. Contar átomos con coeficiente}
\begin{enumerate}
    \item[a)] H: $6$, O: $3$; total $9$.
    \item[b)] C: $2$, O: $4$; total $6$.
    \item[c)] O: $10$; total $10$.
    \item[d)] N: $4$, H: $12$; total $16$.
    \item[e)] Ca: $2$, O: $4$, H: $4$; total $10$.
    \item[f)] Al: $6$, S: $9$, O: $36$; total $51$.
    \item[g)] C: $36$, H: $72$, O: $36$; total $144$.
    \item[h)] H: $10$, N: $10$, O: $30$; total $50$.
\end{enumerate}

\subsubsection*{4. Partículas de un átomo neutro}
\begin{enumerate}
    \item[a)] $p=6$, $n=6$, $e=6$.
    \item[b)] $p=7$, $n=7$, $e=7$.
    \item[c)] $p=8$, $n=8$, $e=8$.
    \item[d)] $p=11$, $n=12$, $e=11$.
    \item[e)] $p=17$, $n=18$, $e=17$.
    \item[f)] $p=17$, $n=20$, $e=17$.
    \item[g)] $p=26$, $n=30$, $e=26$.
    \item[h)] $p=20$, $n=20$, $e=20$.
    \item[i)] $p=13$, $n=14$, $e=13$.
    \item[j)] $p=79$, $n=118$, $e=79$.
\end{enumerate}

\subsubsection*{5. Partículas de un ion}
\begin{enumerate}
    \item[a)] $p=11$, $n=12$, $e=10$.
    \item[b)] $p=17$, $n=18$, $e=18$.
    \item[c)] $p=12$, $n=12$, $e=10$.
    \item[d)] $p=8$, $n=8$, $e=10$.
    \item[e)] $p=13$, $n=14$, $e=10$.
    \item[f)] $p=16$, $n=16$, $e=18$.
    \item[g)] $p=26$, $n=30$, $e=23$.
    \item[h)] $p=20$, $n=20$, $e=18$.
    \item[i)] $p=9$, $n=10$, $e=10$.
    \item[j)] $p=19$, $n=20$, $e=18$.
\end{enumerate}

\subsubsection*{6. De las partículas al símbolo}
\begin{enumerate}
    \item[a)] ${}^{23}\mathrm{Na}^{+}$ ($A=23$, carga $+1$).
    \item[b)] ${}^{35}\mathrm{Cl}^{-}$ ($A=35$, carga $-1$).
    \item[c)] ${}^{16}\mathrm{O}$ ($A=16$, carga $0$).
    \item[d)] ${}^{56}\mathrm{Fe}^{3+}$ ($A=56$, carga $+3$).
    \item[e)] ${}^{32}\mathrm{S}^{2-}$ ($A=32$, carga $-2$).
    \item[f)] ${}^{40}\mathrm{Ca}^{2+}$ ($A=40$, carga $+2$).
    \item[g)] ${}^{27}\mathrm{Al}^{3+}$ ($A=27$, carga $+3$).
    \item[h)] ${}^{80}\mathrm{Br}^{-}$ ($A=80$, carga $-1$).
\end{enumerate}

\subsubsection*{7. Isótopos}
\begin{enumerate}
    \item[a)] $p=1$, $n=0$, $e=1$.
    \item[b)] $p=1$, $n=1$, $e=1$.
    \item[c)] $p=1$, $n=2$, $e=1$.
    \item[d)] $p=6$, $n=6$, $e=6$.
    \item[e)] $p=6$, $n=7$, $e=6$.
    \item[f)] $p=6$, $n=8$, $e=6$.
    \item[g)] $\approx35.48\,\mathrm{u}$.
    \item[h)] $\approx10.80\,\mathrm{u}$.
    \item[i)] $\approx63.62\,\mathrm{u}$.
    \item[j)] $\approx24.32\,\mathrm{u}$.
    \item[k)] $\approx6.92\,\mathrm{u}$.
\end{enumerate}

\subsubsection*{8. Masa molecular}
\begin{enumerate}
    \item[a)] $18\,\mathrm{u}$.
    \item[b)] $44\,\mathrm{u}$.
    \item[c)] $58.5\,\mathrm{u}$.
    \item[d)] $98\,\mathrm{u}$.
    \item[e)] $100\,\mathrm{u}$.
    \item[f)] $17\,\mathrm{u}$.
    \item[g)] $180\,\mathrm{u}$.
    \item[h)] $74\,\mathrm{u}$.
    \item[i)] $342\,\mathrm{u}$.
    \item[j)] $63\,\mathrm{u}$.
    \item[k)] $160\,\mathrm{u}$.
    \item[l)] $159.5\,\mathrm{u}$.
    \item[m)] $158\,\mathrm{u}$.
    \item[n)] $46\,\mathrm{u}$.
    \item[o)] $262\,\mathrm{u}$.
\end{enumerate}

\subsubsection*{9. ¿Está ajustada?}
\begin{enumerate}
    \item[a)] H $4|4$, O $2|2$: sí.
    \item[b)] H $2|2$, O $2|1$: no.
    \item[c)] C $1|1$, H $4|4$, O $4|4$: sí.
    \item[d)] N $2|2$, H $2|6$: no.
    \item[e)] Na $2|2$, Cl $2|2$: sí.
    \item[f)] Fe $2|2$, O $6|4$, C $3|2$: no.
    \item[g)] C $3|3$, H $8|8$, O $10|10$: sí.
    \item[h)] Al $2|2$, O $6|3$: no.
\end{enumerate}

\subsubsection*{10. Ajustar reacciones (nivel 1)}
\begin{enumerate}
    \item[a)] $2\,\mathrm{H_2}+\mathrm{O_2}\rightarrow2\,\mathrm{H_2O}$.
    \item[b)] $\mathrm{N_2}+3\,\mathrm{H_2}\rightarrow2\,\mathrm{NH_3}$.
    \item[c)] $2\,\mathrm{Na}+\mathrm{Cl_2}\rightarrow2\,\mathrm{NaCl}$.
    \item[d)] $2\,\mathrm{Mg}+\mathrm{O_2}\rightarrow2\,\mathrm{MgO}$.
    \item[e)] $\mathrm{H_2}+\mathrm{Cl_2}\rightarrow2\,\mathrm{HCl}$.
    \item[f)] $\mathrm{Fe}+\mathrm{S}\rightarrow\mathrm{FeS}$.
    \item[g)] $4\,\mathrm{K}+\mathrm{O_2}\rightarrow2\,\mathrm{K_2O}$.
    \item[h)] $2\,\mathrm{Al}+3\,\mathrm{Cl_2}\rightarrow2\,\mathrm{AlCl_3}$.
    \item[i)] $\mathrm{CaCO_3}\rightarrow\mathrm{CaO}+\mathrm{CO_2}$.
    \item[j)] $2\,\mathrm{H_2O_2}\rightarrow2\,\mathrm{H_2O}+\mathrm{O_2}$.
    \item[k)] $\mathrm{Zn}+2\,\mathrm{HCl}\rightarrow
    \mathrm{ZnCl_2}+\mathrm{H_2}$.
    \item[l)] $\mathrm{CH_4}+2\,\mathrm{O_2}\rightarrow
    \mathrm{CO_2}+2\,\mathrm{H_2O}$.
\end{enumerate}

\subsubsection*{11. Ajustar reacciones (nivel 2)}
\begin{enumerate}
    \item[a)] $4\,\mathrm{Fe}+3\,\mathrm{O_2}\rightarrow2\,\mathrm{Fe_2O_3}$.
    \item[b)] $4\,\mathrm{Al}+3\,\mathrm{O_2}\rightarrow2\,\mathrm{Al_2O_3}$.
    \item[c)] $2\,\mathrm{C_2H_6}+7\,\mathrm{O_2}\rightarrow
    4\,\mathrm{CO_2}+6\,\mathrm{H_2O}$.
    \item[d)] $\mathrm{C_3H_8}+5\,\mathrm{O_2}\rightarrow
    3\,\mathrm{CO_2}+4\,\mathrm{H_2O}$.
    \item[e)] $2\,\mathrm{Na}+2\,\mathrm{H_2O}\rightarrow
    2\,\mathrm{NaOH}+\mathrm{H_2}$.
    \item[f)] $\mathrm{Mg}+2\,\mathrm{HCl}\rightarrow
    \mathrm{MgCl_2}+\mathrm{H_2}$.
    \item[g)] $\mathrm{Fe_2O_3}+3\,\mathrm{CO}\rightarrow
    2\,\mathrm{Fe}+3\,\mathrm{CO_2}$.
    \item[h)] $2\,\mathrm{KClO_3}\rightarrow2\,\mathrm{KCl}+3\,\mathrm{O_2}$.
    \item[i)] $\mathrm{N_2}+2\,\mathrm{O_2}\rightarrow2\,\mathrm{NO_2}$.
    \item[j)] $4\,\mathrm{P}+5\,\mathrm{O_2}\rightarrow2\,\mathrm{P_2O_5}$.
    \item[k)] $\mathrm{CaCO_3}+2\,\mathrm{HCl}\rightarrow
    \mathrm{CaCl_2}+\mathrm{CO_2}+\mathrm{H_2O}$.
    \item[l)] $2\,\mathrm{Al}+6\,\mathrm{HCl}\rightarrow
    2\,\mathrm{AlCl_3}+3\,\mathrm{H_2}$.
\end{enumerate}

\subsubsection*{12. Ajustar reacciones (nivel 3)}
\begin{enumerate}
    \item[a)] $2\,\mathrm{C_4H_{10}}+13\,\mathrm{O_2}\rightarrow
    8\,\mathrm{CO_2}+10\,\mathrm{H_2O}$.
    \item[b)] $\mathrm{C_2H_5OH}+3\,\mathrm{O_2}\rightarrow
    2\,\mathrm{CO_2}+3\,\mathrm{H_2O}$.
    \item[c)] $\mathrm{C_6H_{12}O_6}+6\,\mathrm{O_2}\rightarrow
    6\,\mathrm{CO_2}+6\,\mathrm{H_2O}$.
    \item[d)] $3\,\mathrm{Fe}+4\,\mathrm{H_2O}\rightarrow
    \mathrm{Fe_3O_4}+4\,\mathrm{H_2}$.
    \item[e)] $4\,\mathrm{NH_3}+5\,\mathrm{O_2}\rightarrow
    4\,\mathrm{NO}+6\,\mathrm{H_2O}$.
    \item[f)] $2\,\mathrm{Al}+3\,\mathrm{H_2SO_4}\rightarrow
    \mathrm{Al_2(SO_4)_3}+3\,\mathrm{H_2}$.
    \item[g)] $3\,\mathrm{Ca(OH)_2}+2\,\mathrm{H_3PO_4}\rightarrow
    \mathrm{Ca_3(PO_4)_2}+6\,\mathrm{H_2O}$.
    \item[h)] $4\,\mathrm{FeS_2}+11\,\mathrm{O_2}\rightarrow
    2\,\mathrm{Fe_2O_3}+8\,\mathrm{SO_2}$.
    \item[i)] $2\,\mathrm{C_8H_{18}}+25\,\mathrm{O_2}\rightarrow
    16\,\mathrm{CO_2}+18\,\mathrm{H_2O}$.
    \item[j)] $2\,\mathrm{Pb(NO_3)_2}\rightarrow
    2\,\mathrm{PbO}+4\,\mathrm{NO_2}+\mathrm{O_2}$.
\end{enumerate}

\subsubsection*{13. Conservación de la masa}
\begin{enumerate}
    \item[a)] $36\,\mathrm{g}$.
    \item[b)] $44\,\mathrm{g}$.
    \item[c)] $44\,\mathrm{g}$.
    \item[d)] $32\,\mathrm{g}$.
    \item[e)] $1.6\,\mathrm{g}$.
    \item[f)] $6.5\,\mathrm{g}$.
    \item[g)] $36\,\mathrm{g}$.
    \item[h)] $1\,\mathrm{g}$.
\end{enumerate}

\subsubsection*{14. Proporciones fijas en masa}
\begin{enumerate}
    \item[a)] $24\,\mathrm{g}$.
    \item[b)] $5\,\mathrm{g}$.
    \item[c)] $18\,\mathrm{g}$.
    \item[d)] $16\,\mathrm{g}$.
    \item[e)] $30\,\mathrm{g}$.
    \item[f)] $16\,\mathrm{g}$.
    \item[g)] $35\,\mathrm{g}$ de Fe necesitan $20\,\mathrm{g}$ de S;
    sobran $15\,\mathrm{g}$ de S;
    se forman $55\,\mathrm{g}$ de $\mathrm{FeS}$.
\end{enumerate}

\subsubsection*{15. Moles y masa}
\begin{enumerate}
    \item[a)] $2\,\mathrm{mol}$.
    \item[b)] $2\,\mathrm{mol}$.
    \item[c)] $29.25\,\mathrm{g}$.
    \item[d)] $96\,\mathrm{g}$.
    \item[e)] $0.05\,\mathrm{mol}$.
    \item[f)] $25\,\mathrm{g}$.
    \item[g)] $0.25\,\mathrm{mol}$.
    \item[h)] $270\,\mathrm{g}$.
    \item[i)] $0.5\,\mathrm{mol}$.
    \item[j)] $0.05\,\mathrm{mol}$.
    \item[k)] $31.6\,\mathrm{g}$.
    \item[l)] $10\,\mathrm{mol}$.
\end{enumerate}

\subsubsection*{16. Moles y partículas}
\begin{enumerate}
    \item[a)] $1.2044\cdot10^{24}$ moléculas.
    \item[b)] $3.011\cdot10^{23}$ átomos.
    \item[c)] $1\,\mathrm{mol}$;
    $6.022\cdot10^{23}$ moléculas;
    $1.2044\cdot10^{24}$ átomos de H.
    \item[d)] $1\,\mathrm{mol}$;
    $6.022\cdot10^{23}$ moléculas;
    $1.2044\cdot10^{24}$ átomos de O.
    \item[e)] $2\,\mathrm{mol}$; $64\,\mathrm{g}$.
    \item[f)] $0.5\,\mathrm{mol}$; $11.5\,\mathrm{g}$.
    \item[g)] $1.8066\cdot10^{24}$ átomos de H.
    \item[h)] $0.2\,\mathrm{mol}$;
    $1.2044\cdot10^{23}$ moléculas.
\end{enumerate}

\subsubsection*{17. Moles en una reacción ajustada}
\begin{enumerate}
    \item[a)] $6\,\mathrm{mol}$ de $\mathrm{H_2}$,
    $4\,\mathrm{mol}$ de $\mathrm{NH_3}$.
    \item[b)] $3\,\mathrm{mol}$ de $\mathrm{N_2}$,
    $6\,\mathrm{mol}$ de $\mathrm{NH_3}$.
    \item[c)] $5\,\mathrm{mol}$ de $\mathrm{N_2}$,
    $15\,\mathrm{mol}$ de $\mathrm{H_2}$.
    \item[d)] $6\,\mathrm{mol}$ de $\mathrm{O_2}$,
    $3\,\mathrm{mol}$ de $\mathrm{CO_2}$,
    $6\,\mathrm{mol}$ de $\mathrm{H_2O}$.
    \item[e)] $2.5\,\mathrm{mol}$ de $\mathrm{CH_4}$,
    $5\,\mathrm{mol}$ de $\mathrm{H_2O}$.
    \item[f)] $6\,\mathrm{mol}$ de $\mathrm{Cl_2}$,
    $4\,\mathrm{mol}$ de $\mathrm{AlCl_3}$.
    \item[g)] $0.4\,\mathrm{mol}$ de Al,
    $0.4\,\mathrm{mol}$ de $\mathrm{AlCl_3}$.
    \item[h)] $10\,\mathrm{mol}$ de $\mathrm{O_2}$,
    $6\,\mathrm{mol}$ de $\mathrm{CO_2}$,
    $8\,\mathrm{mol}$ de $\mathrm{H_2O}$.
\end{enumerate}

\subsubsection*{18. Cálculos con masas}
\begin{enumerate}
    \item[a)] $2\,\mathrm{H_2}+\mathrm{O_2}\rightarrow2\,\mathrm{H_2O}$;
    $n(\mathrm{H_2})=4\,\mathrm{mol}$;
    $m(\mathrm{H_2O})=72\,\mathrm{g}$.

    \item[b)] $\mathrm{CaCO_3}\rightarrow\mathrm{CaO}+\mathrm{CO_2}$;
    $n(\mathrm{CaCO_3})=0.5\,\mathrm{mol}$;
    $m(\mathrm{CaO})=28\,\mathrm{g}$.

    \item[c)] $\mathrm{CH_4}+2\,\mathrm{O_2}\rightarrow
    \mathrm{CO_2}+2\,\mathrm{H_2O}$;
    $n(\mathrm{CH_4})=2\,\mathrm{mol}$;
    $m(\mathrm{O_2})=128\,\mathrm{g}$.

    \item[d)] $\mathrm{CH_4}+2\,\mathrm{O_2}\rightarrow
    \mathrm{CO_2}+2\,\mathrm{H_2O}$;
    $n(\mathrm{CH_4})=2\,\mathrm{mol}$;
    $m(\mathrm{CO_2})=88\,\mathrm{g}$.

    \item[e)] $2\,\mathrm{Mg}+\mathrm{O_2}\rightarrow2\,\mathrm{MgO}$;
    $n(\mathrm{Mg})=0.5\,\mathrm{mol}$;
    $m(\mathrm{MgO})=20\,\mathrm{g}$.

    \item[f)] $\mathrm{N_2}+3\,\mathrm{H_2}\rightarrow2\,\mathrm{NH_3}$;
    $n(\mathrm{N_2})=2\,\mathrm{mol}$;
    $m(\mathrm{NH_3})=68\,\mathrm{g}$.

    \item[g)] $\mathrm{Fe_2O_3}+3\,\mathrm{CO}\rightarrow
    2\,\mathrm{Fe}+3\,\mathrm{CO_2}$;
    $n(\mathrm{Fe_2O_3})=2\,\mathrm{mol}$;
    $m(\mathrm{Fe})=224\,\mathrm{g}$.

    \item[h)] $2\,\mathrm{KClO_3}\rightarrow2\,\mathrm{KCl}+3\,\mathrm{O_2}$;
    $n(\mathrm{KClO_3})=0.2\,\mathrm{mol}$;
    $m(\mathrm{O_2})=9.6\,\mathrm{g}$.

    \item[i)] $\mathrm{Zn}+2\,\mathrm{HCl}\rightarrow
    \mathrm{ZnCl_2}+\mathrm{H_2}$;
    $n(\mathrm{Zn})=0.2\,\mathrm{mol}$;
    $m(\mathrm{H_2})=0.4\,\mathrm{g}$.
\end{enumerate}

\subsubsection*{19. Volumen de gases en condiciones normales}
\begin{enumerate}
    \item[a)] $44.8\,\mathrm{L}$.
    \item[b)] $0.5\,\mathrm{mol}$; $22\,\mathrm{g}$.
    \item[c)] $44.8\,\mathrm{L}$.
    \item[d)] $7\,\mathrm{g}$.
    \item[e)] $11.2\,\mathrm{L}$.
    \item[f)] $3\,\mathrm{mol}$; $51\,\mathrm{g}$.
    \item[g)] $2.24\,\mathrm{L}$.
    \item[h)] $11.2\,\mathrm{L}$.
\end{enumerate}

\subsubsection*{20. Volúmenes en las reacciones}
\begin{enumerate}
    \item[a)] $30\,\mathrm{L}$ de $\mathrm{H_2}$,
    $20\,\mathrm{L}$ de $\mathrm{NH_3}$.
    \item[b)] $3\,\mathrm{L}$ de $\mathrm{O_2}$,
    $6\,\mathrm{L}$ de $\mathrm{H_2O}$.
    \item[c)] $10\,\mathrm{L}$ de $\mathrm{O_2}$,
    $5\,\mathrm{L}$ de $\mathrm{CO_2}$,
    $10\,\mathrm{L}$ de $\mathrm{H_2O}$.
    \item[d)] $15\,\mathrm{L}$ de $\mathrm{H_2}$,
    $30\,\mathrm{L}$ de $\mathrm{HCl}$.
    \item[e)] $1\,\mathrm{mol}$ de $\mathrm{CO}\rightarrow
    0.5\,\mathrm{mol}$ de $\mathrm{O_2}\rightarrow11.2\,\mathrm{L}$.
    \item[f)] $0.5\,\mathrm{mol}\rightarrow11.2\,\mathrm{L}$.
    \item[g)] $0.1\,\mathrm{mol}\rightarrow2.24\,\mathrm{L}$.
    \item[h)] $0.5\,\mathrm{mol}$ de $\mathrm{C_3H_8}\rightarrow
    2.5\,\mathrm{mol}$ de $\mathrm{O_2}=56\,\mathrm{L}$;
    $1.5\,\mathrm{mol}$ de $\mathrm{CO_2}=33.6\,\mathrm{L}$.
\end{enumerate}

\subsection{Soluciones de 4.º ESO}

\subsubsection*{21. Configuración electrónica}
\begin{enumerate}
    \item[a)] $1\mathrm{s}^1$.
    \item[b)] $1\mathrm{s}^2\ 2\mathrm{s}^2\ 2\mathrm{p}^2$.
    \item[c)] $1\mathrm{s}^2\ 2\mathrm{s}^2\ 2\mathrm{p}^3$.
    \item[d)] $1\mathrm{s}^2\ 2\mathrm{s}^2\ 2\mathrm{p}^4$.
    \item[e)] $1\mathrm{s}^2\ 2\mathrm{s}^2\ 2\mathrm{p}^5$.
    \item[f)] $1\mathrm{s}^2\ 2\mathrm{s}^2\ 2\mathrm{p}^6$.
    \item[g)] $1\mathrm{s}^2\ 2\mathrm{s}^2\ 2\mathrm{p}^6\ 3\mathrm{s}^1$.
    \item[h)] $1\mathrm{s}^2\ 2\mathrm{s}^2\ 2\mathrm{p}^6\
    3\mathrm{s}^2\ 3\mathrm{p}^1$.
    \item[i)] $1\mathrm{s}^2\ 2\mathrm{s}^2\ 2\mathrm{p}^6\
    3\mathrm{s}^2\ 3\mathrm{p}^4$.
    \item[j)] $1\mathrm{s}^2\ 2\mathrm{s}^2\ 2\mathrm{p}^6\
    3\mathrm{s}^2\ 3\mathrm{p}^5$.
    \item[k)] $1\mathrm{s}^2\ 2\mathrm{s}^2\ 2\mathrm{p}^6\
    3\mathrm{s}^2\ 3\mathrm{p}^6\ 4\mathrm{s}^1$.
    \item[l)] $1\mathrm{s}^2\ 2\mathrm{s}^2\ 2\mathrm{p}^6\
    3\mathrm{s}^2\ 3\mathrm{p}^6\ 4\mathrm{s}^2$.
    \item[m)] $1\mathrm{s}^2\ 2\mathrm{s}^2\ 2\mathrm{p}^6\
    3\mathrm{s}^2\ 3\mathrm{p}^6\ 4\mathrm{s}^2\ 3\mathrm{d}^6$.
    \item[n)] $1\mathrm{s}^2\ 2\mathrm{s}^2\ 2\mathrm{p}^6\
    3\mathrm{s}^2\ 3\mathrm{p}^6\ 4\mathrm{s}^2\
    3\mathrm{d}^{10}\ 4\mathrm{p}^5$.
\end{enumerate}

\subsubsection*{22. Grupo y periodo}
\begin{enumerate}
    \item[a)] $1\mathrm{s}^2\ 2\mathrm{s}^1$;
    grupo $1$, periodo $2$, $1\,e^{-}$ de valencia.

    \item[b)] $1\mathrm{s}^2\ 2\mathrm{s}^2\ 2\mathrm{p}^3$;
    grupo $15$, periodo $2$, $5\,e^{-}$ de valencia.

    \item[c)] $1\mathrm{s}^2\ 2\mathrm{s}^2\ 2\mathrm{p}^5$;
    grupo $17$, periodo $2$, $7\,e^{-}$ de valencia.

    \item[d)] $1\mathrm{s}^2\ 2\mathrm{s}^2\ 2\mathrm{p}^6\ 3\mathrm{s}^2$;
    grupo $2$, periodo $3$, $2\,e^{-}$ de valencia.

    \item[e)] $1\mathrm{s}^2\ 2\mathrm{s}^2\ 2\mathrm{p}^6\
    3\mathrm{s}^2\ 3\mathrm{p}^1$;
    grupo $13$, periodo $3$, $3\,e^{-}$ de valencia.

    \item[f)] $1\mathrm{s}^2\ 2\mathrm{s}^2\ 2\mathrm{p}^6\
    3\mathrm{s}^2\ 3\mathrm{p}^3$;
    grupo $15$, periodo $3$, $5\,e^{-}$ de valencia.

    \item[g)] $1\mathrm{s}^2\ 2\mathrm{s}^2\ 2\mathrm{p}^6\
    3\mathrm{s}^2\ 3\mathrm{p}^4$;
    grupo $16$, periodo $3$, $6\,e^{-}$ de valencia.

    \item[h)] $1\mathrm{s}^2\ 2\mathrm{s}^2\ 2\mathrm{p}^6\
    3\mathrm{s}^2\ 3\mathrm{p}^5$;
    grupo $17$, periodo $3$, $7\,e^{-}$ de valencia.

    \item[i)] $1\mathrm{s}^2\ 2\mathrm{s}^2\ 2\mathrm{p}^6\
    3\mathrm{s}^2\ 3\mathrm{p}^6$;
    grupo $18$, periodo $3$, $8\,e^{-}$ de valencia.

    \item[j)] $1\mathrm{s}^2\ 2\mathrm{s}^2\ 2\mathrm{p}^6\
    3\mathrm{s}^2\ 3\mathrm{p}^6\ 4\mathrm{s}^1$;
    grupo $1$, periodo $4$, $1\,e^{-}$ de valencia.

    \item[k)] $1\mathrm{s}^2\ 2\mathrm{s}^2\ 2\mathrm{p}^6\
    3\mathrm{s}^2\ 3\mathrm{p}^6\ 4\mathrm{s}^2$;
    grupo $2$, periodo $4$, $2\,e^{-}$ de valencia.

    \item[l)] $1\mathrm{s}^2\ 2\mathrm{s}^2\ 2\mathrm{p}^6\
    3\mathrm{s}^2\ 3\mathrm{p}^6\ 4\mathrm{s}^2\
    3\mathrm{d}^{10}\ 4\mathrm{p}^5$;
    grupo $17$, periodo $4$, $7\,e^{-}$ de valencia.
\end{enumerate}

\subsubsection*{23. Ion más probable}
\begin{enumerate}
    \item[a)] Grupo $1$; $\mathrm{Li}^{+}$;
    $2$ electrones, como el He.
    \item[b)] Grupo $2$; $\mathrm{Mg}^{2+}$;
    $10$ electrones, como el Ne.
    \item[c)] Grupo $13$; $\mathrm{Al}^{3+}$;
    $10$ electrones, como el Ne.
    \item[d)] Grupo $15$; $\mathrm{N}^{3-}$;
    $10$ electrones, como el Ne.
    \item[e)] Grupo $16$; $\mathrm{O}^{2-}$;
    $10$ electrones, como el Ne.
    \item[f)] Grupo $17$; $\mathrm{F}^{-}$;
    $10$ electrones, como el Ne.
    \item[g)] Grupo $1$; $\mathrm{Na}^{+}$;
    $10$ electrones, como el Ne.
    \item[h)] Grupo $16$; $\mathrm{S}^{2-}$;
    $18$ electrones, como el Ar.
    \item[i)] Grupo $17$; $\mathrm{Cl}^{-}$;
    $18$ electrones, como el Ar.
    \item[j)] Grupo $2$; $\mathrm{Ca}^{2+}$;
    $18$ electrones, como el Ar.
    \item[k)] Grupo $1$; $\mathrm{K}^{+}$;
    $18$ electrones, como el Ar.
    \item[l)] Grupo $17$; $\mathrm{Br}^{-}$;
    $36$ electrones, como el Kr.
\end{enumerate}

\subsubsection*{24. Fórmulas de compuestos iónicos}
\begin{enumerate}
    \item[a)] $\mathrm{NaCl}$.
    \item[b)] $\mathrm{MgCl_2}$.
    \item[c)] $\mathrm{Al_2O_3}$.
    \item[d)] $\mathrm{CaO}$.
    \item[e)] $\mathrm{K_2S}$.
    \item[f)] $\mathrm{AlCl_3}$.
    \item[g)] $\mathrm{Fe_2O_3}$.
    \item[h)] $\mathrm{Mg_3N_2}$.
    \item[i)] $\mathrm{Li_2O}$.
    \item[j)] $\mathrm{CaF_2}$.
    \item[k)] $\mathrm{Al_2S_3}$.
    \item[l)] $\mathrm{BaBr_2}$.
\end{enumerate}

\subsubsection*{25. Ajustar reacciones (nivel 4)}
\begin{enumerate}
    \item[a)] $\mathrm{C_5H_{12}}+8\,\mathrm{O_2}\rightarrow
    5\,\mathrm{CO_2}+6\,\mathrm{H_2O}$.
    \item[b)] $2\,\mathrm{C_3H_8O_3}+7\,\mathrm{O_2}\rightarrow
    6\,\mathrm{CO_2}+8\,\mathrm{H_2O}$.
    \item[c)] $\mathrm{C_{12}H_{22}O_{11}}+12\,\mathrm{O_2}\rightarrow
    12\,\mathrm{CO_2}+11\,\mathrm{H_2O}$.
    \item[d)] $2\,\mathrm{Fe_2O_3}+3\,\mathrm{C}\rightarrow
    4\,\mathrm{Fe}+3\,\mathrm{CO_2}$.
    \item[e)] $2\,\mathrm{Al}+\mathrm{Fe_2O_3}\rightarrow
    \mathrm{Al_2O_3}+2\,\mathrm{Fe}$.
    \item[f)] $\mathrm{NH_4NO_3}\rightarrow\mathrm{N_2O}+2\,\mathrm{H_2O}$.
    \item[g)] $\mathrm{MnO_2}+4\,\mathrm{HCl}\rightarrow
    \mathrm{MnCl_2}+\mathrm{Cl_2}+2\,\mathrm{H_2O}$.
    \item[h)] $3\,\mathrm{Cu}+8\,\mathrm{HNO_3}\rightarrow
    3\,\mathrm{Cu(NO_3)_2}+2\,\mathrm{NO}+4\,\mathrm{H_2O}$.
    \item[i)] $2\,\mathrm{KMnO_4}+16\,\mathrm{HCl}\rightarrow
    2\,\mathrm{KCl}+2\,\mathrm{MnCl_2}+5\,\mathrm{Cl_2}+8\,\mathrm{H_2O}$.
    \item[j)] $2\,\mathrm{Ca_3(PO_4)_2}+6\,\mathrm{SiO_2}
    +10\,\mathrm{C}\rightarrow
    6\,\mathrm{CaSiO_3}+\mathrm{P_4}+10\,\mathrm{CO}$.
\end{enumerate}

\subsubsection*{26. Cálculos completos}
\begin{enumerate}
    \item[a)] $2\,\mathrm{mol}$ de $\mathrm{H_2O_2}\rightarrow
    1\,\mathrm{mol}$ de $\mathrm{O_2}=22.4\,\mathrm{L}$.

    \item[b)] $0.5\,\mathrm{mol}$ de Zn $\rightarrow
    0.5\,\mathrm{mol}$ de $\mathrm{H_2}=11.2\,\mathrm{L}$.

    \item[c)] $3\,\mathrm{mol}$ de $\mathrm{CO_2}\rightarrow
    1\,\mathrm{mol}$ de $\mathrm{C_3H_8}=44\,\mathrm{g}$.

    \item[d)] $0.2\,\mathrm{mol}$ de Na $\rightarrow
    0.2\,\mathrm{mol}$ de $\mathrm{NaOH}=8\,\mathrm{g}$;
    $0.1\,\mathrm{mol}$ de $\mathrm{H_2}=2.24\,\mathrm{L}$.

    \item[e)] $0.25\,\mathrm{mol}$ de $\mathrm{CaCO_3}\rightarrow
    27.75\,\mathrm{g}$ de $\mathrm{CaCl_2}$;
    $5.6\,\mathrm{L}$ de $\mathrm{CO_2}$.

    \item[f)] $0.2\,\mathrm{mol}$ de Fe $\rightarrow
    0.1\,\mathrm{mol}$ de $\mathrm{Fe_2O_3}=16\,\mathrm{g}$;
    $0.15\,\mathrm{mol}$ de $\mathrm{O_2}=3.36\,\mathrm{L}$.

    \item[g)] $2\,\mathrm{mol}$ de $\mathrm{H_2}\rightarrow
    \frac{4}{3}\,\mathrm{mol}$ de $\mathrm{NH_3}\approx22.67\,\mathrm{g}$.

    \item[h)] $0.5\,\mathrm{mol}$ de glucosa $\rightarrow
    3\,\mathrm{mol}$ de $\mathrm{CO_2}=1.8066\cdot10^{24}$ moléculas.
\end{enumerate}

\subsubsection*{27. Reactivo limitante}
\begin{enumerate}
    \item[a)] $2\,\mathrm{H_2}+\mathrm{O_2}\rightarrow2\,\mathrm{H_2O}$;
    limitante $\mathrm{H_2}$;
    $36\,\mathrm{g}$ de $\mathrm{H_2O}$;
    sobran $8\,\mathrm{g}$ de $\mathrm{O_2}$.

    \item[b)] $\mathrm{N_2}+3\,\mathrm{H_2}\rightarrow2\,\mathrm{NH_3}$;
    limitante $\mathrm{N_2}$;
    $34\,\mathrm{g}$ de $\mathrm{NH_3}$;
    sobran $3\,\mathrm{g}$ de $\mathrm{H_2}$.

    \item[c)] $\mathrm{Zn}+2\,\mathrm{HCl}\rightarrow
    \mathrm{ZnCl_2}+\mathrm{H_2}$;
    limitante Zn;
    $0.4\,\mathrm{g}$ de $\mathrm{H_2}$;
    sobran $3.65\,\mathrm{g}$ de $\mathrm{HCl}$.

    \item[d)] $2\,\mathrm{Al}+3\,\mathrm{Cl_2}\rightarrow2\,\mathrm{AlCl_3}$;
    no sobra ninguno;
    $26.7\,\mathrm{g}$ de $\mathrm{AlCl_3}$.

    \item[e)] $\mathrm{CH_4}+2\,\mathrm{O_2}\rightarrow
    \mathrm{CO_2}+2\,\mathrm{H_2O}$;
    limitante $\mathrm{CH_4}$;
    $22\,\mathrm{g}$ de $\mathrm{CO_2}$;
    sobran $16\,\mathrm{g}$ de $\mathrm{O_2}$.

    \item[f)] $\mathrm{Fe}+\mathrm{S}\rightarrow\mathrm{FeS}$;
    limitante Fe;
    $44\,\mathrm{g}$ de $\mathrm{FeS}$;
    sobran $4\,\mathrm{g}$ de S.
\end{enumerate}

\subsubsection*{28. Rendimiento y pureza}
\begin{enumerate}
    \item[a)] Teórico $56\,\mathrm{g}$; $90\,\%$.

    \item[b)] Teórico $34\,\mathrm{g}$; $75\,\%$.

    \item[c)] Teórico $72\,\mathrm{g}$; real $57.6\,\mathrm{g}$.

    \item[d)] Teórico $0.5\,\mathrm{mol}$;
    real $0.3\,\mathrm{mol}=6.72\,\mathrm{L}$.

    \item[e)] $150\,\mathrm{g}$ de $\mathrm{CaCO_3}
    =1.5\,\mathrm{mol}\rightarrow66\,\mathrm{g}$ de $\mathrm{CO_2}$.

    \item[f)] $13.08\,\mathrm{g}$ de Zn $=0.2\,\mathrm{mol}
    \rightarrow4.48\,\mathrm{L}$.

    \item[g)] $0.75\,\mathrm{mol}$ de Fe $\rightarrow
    0.375\,\mathrm{mol}$ de $\mathrm{Fe_2O_3}=60\,\mathrm{g}$;
    $75\,\%$.

    \item[h)] $90\,\mathrm{g}$ de $\mathrm{CaCO_3}
    =0.9\,\mathrm{mol}\rightarrow$
    teórico $50.4\,\mathrm{g}$; real $40.32\,\mathrm{g}$.
\end{enumerate}

\subsubsection*{29. Disoluciones}
\begin{enumerate}
    \item[a)] $40\,\mathrm{g/L}$.
    \item[b)] $2\,\mathrm{M}$.
    \item[c)] $0.1\,\mathrm{mol}\rightarrow0.5\,\mathrm{M}$.
    \item[d)] $0.1\,\mathrm{mol}\rightarrow0.2\,\mathrm{M}$.
    \item[e)] $0.05\,\mathrm{mol}\rightarrow2.925\,\mathrm{g}$.
    \item[f)] $150\,\mathrm{g}$ de disolución $\rightarrow10\,\%$.
    \item[g)] $15\,\mathrm{g}$.
    \item[h)] $200\,\mathrm{mL}$.
    \item[i)] $0.1\,\mathrm{mol}$ en $0.25\,\mathrm{L}\rightarrow0.4\,\mathrm{M}$.
    \item[j)] $29.25\,\mathrm{g/L}$.
    \item[k)] $0.2\,\mathrm{mol}\rightarrow36\,\mathrm{g}$.
    \item[l)] $1\,\mathrm{M}$.
\end{enumerate}

\subsubsection*{30. Reacciones en disolución}
\begin{enumerate}
    \item[a)] $0.1\,\mathrm{mol}$ de Zn $\rightarrow
    0.2\,\mathrm{mol}$ de $\mathrm{HCl}\rightarrow100\,\mathrm{mL}$.

    \item[b)] $0.005\,\mathrm{mol}$ de $\mathrm{NaOH}\rightarrow
    0.005\,\mathrm{mol}$ de $\mathrm{HCl}\rightarrow50\,\mathrm{mL}$.

    \item[c)] $0.01\,\mathrm{mol}$ de $\mathrm{H_2SO_4}\rightarrow
    0.02\,\mathrm{mol}$ de $\mathrm{NaOH}\rightarrow20\,\mathrm{mL}$.

    \item[d)] $0.1\,\mathrm{mol}$ de $\mathrm{HCl}\rightarrow
    0.05\,\mathrm{mol}$ de $\mathrm{CaCO_3}=5\,\mathrm{g}$;
    $0.05\,\mathrm{mol}$ de $\mathrm{CO_2}=1.12\,\mathrm{L}$.

    \item[e)] Mg $0.1\,\mathrm{mol}$,
    $\mathrm{HCl}$ $0.1\,\mathrm{mol}$:
    limitante $\mathrm{HCl}$;
    $0.05\,\mathrm{mol}$ de $\mathrm{H_2}=1.12\,\mathrm{L}$;
    sobran $0.05\,\mathrm{mol}$ de Mg $=1.2\,\mathrm{g}$.

    \item[f)] $0.005\,\mathrm{mol}\rightarrow0.7175\,\mathrm{g}$.
\end{enumerate}